\section{A larger shrinking family with a finite-count envelope}
\label{sec:v88-envelope}
The moment in Theorem~\ref{thm:v88-moment} improves the range of
initial geometries that can be controlled before the collision limit.
We write the trace argument for the full moment range, so its use does
not rely on substituting an exponent into the narrower historical
joint-moment theorem.

Fix $1<p<3/2$ and put $e_p=1-1/p$. At roof $t$ select every actual
maximal segment satisfying
\[
 a_j<t<b_j,\qquad z_j(t)\ge a,\qquad b_j-a_j\ge\delta.
\]
Denote this Borel family by $\mathcal E_t(a,\delta)$; $a$ is a
gradient threshold, not the endpoint roof $a_j$. Put
\[
 d=\min\{\delta/16,a^2/(32\Lambda_\circ),1/16\},\qquad
 J_t=(t-4d,t+4d),\qquad M_p(J)=\int_Ju_p.
\]
There is no restriction on internal incidences or on $I_m$.

\begin{lemma}[Complete traces for the enlarged moment range]
\label{lem:v88-trace}
For $0<a,\delta\le1$, in the common essential-roof versions,
\begin{equation}\label{eq:v88-trace}
 S_t:=\sum_j\int_{\mathcal E_t(a,\delta)}
    \sup_{t\le v<\min(t+d,b_j)}A_j(v,\theta)\,d\theta
                      \le C_p d^{-1}M_p(J_t).
\end{equation}
The displayed sum is over the disjoint segment parameter space.
\end{lemma}
\begin{proof}
One side of a selected segment has length at least $\delta/2\ge8d$.
Choose the anchor interval $(t-3d,t-2d)$ if the backward side is
long, and $(t+2d,t+3d)$ otherwise, with that Borel tie rule.
Every joining path from the anchor to a target stays in the same
physical segment and in $J_t$. Radial growth gives
$z(v)^2\ge a^2-8\Lambda_\circ d\ge3a^2/4$ on these paths.
The exact angular density therefore satisfies $|A'|\le(C/a)AI_m$.

For $f=A^{1/p}$ one has $|f'|\le(C/(pa))fI_m$ locally. Choose
an anchor point with $f\le(d^{-1}\int_{\rm anchor}A)^{1/p}$,
allowing an arbitrarily small error. The fundamental theorem of
calculus and H\"older on a joining interval of length at most $8d$
give
\[
 \sup f\le(d^{-1}\!\int_{J_t\cap(a_j,b_j)}A)^{1/p}
       +\frac C{pa}(8d)^{1-1/p}
                  (\int_{J_t\cap(a_j,b_j)}AI_m^p)^{1/p}.
\]
Since $d\le a$, raising to the power $p$ bounds the trace by
$C_pd^{-1}\int AI_m^p$. The Borel segment bijection identifies
the angular sum of this energy with a submeasure of $M_p(J_t)$,
with multiplicity one. Countable dense roof offsets give measurable
suprema; positive exhaustion gives the endpoint limits. This proves
\eqref{eq:v88-trace} for every $p>1$ for which the energy is finite.
\end{proof}

For $0<h\le d$ define the complete reachable variation source
\[
 \mathcal V_h^{a,\delta}(t)=\sum_j\int_{\mathcal E_t(a,\delta)}
    \int_0^{\min(h,b_j-t)}
       A_j(t+s,\theta)I_m(t+s,\theta)/z_j(t+s,\theta)
                                                     \,ds\,d\theta.
\]
\begin{theorem}[Trace-to-variation bound at every finite count]
\label{thm:v88-variation}
For every such $a,\delta,h$,
\begin{equation}\label{eq:v88-variation}
 \mathcal V_h^{a,\delta}(t)
 \le \frac{h^{e_p}}a S_t^{e_p}M_p((t,t+h))^{1/p}
 \le \frac{C_p}a(h/d)^{e_p}M_p(J_t).
\end{equation}
The complete sum includes all actual words of the fixed exact label.
\end{theorem}
\begin{proof}
On every selected forward segment $z(t+s)\ge a$. Apply H\"older
first in $s$ to $AI_m=A^{1-1/p}(AI_m^p)^{1/p}$ and then on the
disjoint angular parameter space. The summed first factor is exactly
$S_t$, not a maximum over words. This gives the first bound. Use
Lemma~\ref{lem:v88-trace} and $(t,t+h)\subset J_t$ for the second.
The integrals stop at the first true exit; no later re-entry is used.
\end{proof}

\begin{theorem}[An expanded finite-count microscopic family]
\label{thm:v88-expanded}
Set
\[
 p=10/7,\quad a(h)=h^{1/8},\quad\delta(h)=h^{1/4},\quad
 c_d=\min\{1/16,(32\Lambda_\circ)^{-1}\}.
\]
There is a fixed $h_*>0$ such that, for $0<h<h_*$,
$d=c_dh^{1/4}\ge h$, and
\begin{equation}\label{eq:v88-expanded}
 \sup_{m\ge2,R,n,k}m^2\|\mathcal V_h^{a(h),\delta(h)}\|_\infty
                              \le K_{B_\dagger}h^{1/10}.
\end{equation}
The same bound, up to a physical constant, pays surviving inverse-twist
loss and the original starting mass of genuine grazing first exits
on this family. The estimate is uniform at finite collision count.
\end{theorem}
\begin{proof}
Choose $h_*<\min\{1,c_d^{4/3}\}$. Insert
\eqref{eq:v88-uniform-energy} into \eqref{eq:v88-variation}.
The exponent is
\[
 -\frac18+\left(1-\frac7{10}\right)\left(1-\frac14\right)
                       =\frac1{10}.
\]
The constant is for the same single band at every $m$ and $h$.
On survivors negative optical variation pays the decrease of $\omega$;
at a genuine grazing exit it pays the full initial optical value
because the terminal value is zero. The angular normalization ratios
are uniformly bounded, so the same variation source bounds both
contributions on their disjoint initial sets. This does not assign
zero starting mass to a grazing exit.
\end{proof}

More generally, for $a=h^\xi$ and $\delta=h^\upsilon$, put
$b=\max\{2\xi,\upsilon\}$. The same theorem holds with exponent
$e_p(1-b)-\xi$ whenever $b<1$ and this exponent is positive.
For $\upsilon=2\xi$, every $\xi<1/5$ is admissible by a fixed
choice of $p<3/2$. This is the range delivered by this argument,
not a sharpness assertion for the billiard. The example above includes
states below both of the earlier thresholds $h^{1/128},h^{1/64}$.
It still leaves the complementary original starting source to be
estimated; cofinal cutoffs alone never prove its charge zero.
